% ------------------------------------------------------------------
%  Capitolo 20 — Il semestre, le lezioni, gli esami
%  Parte V
%  Ricerca di riferimento: ricerca 04 §6–7; 01 §5.2; 02 §9
%  Voci bibliografiche principali: rawson2013, rawson2022, kornell2009, freeman2014, meyer2013, bastani2025
%  Epigrafe: ✔ verificata sul testo (vedi ricerca 03 e registro, ricerca 00)
%  Scheda del capitolo: ricerca/06_indice-ragionato.md, capitolo 20
% ------------------------------------------------------------------
\chapter{Il semestre, le lezioni, gli esami}
\label{cap:semestre}

\epigrafe[invecchiare imparando ogni giorno molte cose nuove]{senescere se multa in dies addiscentem}{Cicerone, Cato maior de senectute 26}

\iniziale{Q}{uesto capitolo} \segnaposto{apertura: da scrivere — vedi la scheda nell'indice ragionato}.

\section{Pianificare il semestre, non la sessione}

\segnaposto{da scrivere}

\section{Dalla lezione agli appunti, dagli appunti alle domande}

\segnaposto{da scrivere}

\section{Prepararsi agli esami scritti e orali}

\segnaposto{da scrivere}

\section{Studiare in gruppo, e con l'intelligenza artificiale}

\segnaposto{da scrivere}

\section{Dopo la laurea: imparare per tutta la vita}

\segnaposto{da scrivere}

\begin{daricordare}
\segnaposto{il principio del capitolo in due righe}
\end{daricordare}

\begin{domande}
  \item \segnaposto{domanda di recupero su questo capitolo}
  \item \segnaposto{domanda di applicazione a una materia che studi}
  \item \segnaposto{domanda di ripresa su un capitolo precedente}
\end{domande}

\begin{sintesi}
  \item \segnaposto{punto di sintesi}
\end{sintesi}

\finecapitolo
